% =============================================================================
% Chapter 4: Data Engineering — ML Systems Lecture Slides
% =============================================================================
% IMAGES NEEDED (generate SVGs, convert to PDF):
%   - ch04-dataset-compiler.pdf   - ch04-energy-movement.pdf
%   - ch04-feeding-problem.pdf    - ch04-four-pillars.pdf
%   - ch04-drift-types.pdf        - ch04-training-serving-skew.pdf
%   - ch04-storage-hierarchy.pdf  - ch04-data-debt-categories.pdf
%   - cover_data_engineering.png
% =============================================================================
\documentclass[aspectratio=169, 12pt]{beamer}
\usepackage{../../assets/beamerthememlsys}

\mlsyssetup{
  volume   = {Volume I},
  chapter  = {Chapter 4},
  logo = {../../assets/img/logo-mlsysbook.png},
  instlogo = {../../assets/img/logo-harvard.png},
  chaptertitle = {Data Engineering},
}

% --- Fonts ---
\usepackage{fontspec}
\setsansfont{Helvetica Neue}[
  BoldFont={Helvetica Neue Bold},
  ItalicFont={Helvetica Neue Italic},
  BoldItalicFont={Helvetica Neue Bold Italic},
]
\setmonofont{JetBrains Mono}[Scale=0.85]

% --- Packages ---
\usepackage{booktabs}
\usepackage{amsmath}

% --- Image paths ---
\graphicspath{
  {images/}
}

% --- Chapter-specific macros ---
\newcommand{\DAM}{D\raisebox{0.08em}{\tiny$\bullet$}A\raisebox{0.08em}{\tiny$\bullet$}M}

% --- Helper: safe image include ---
\newcommand{\safeimg}[2][width=\textwidth,keepaspectratio]{%
  \IfFileExists{images/#2}{\includegraphics[#1]{#2}}{%
    \IfFileExists{#2}{\includegraphics[#1]{#2}}{%
      \fbox{\parbox[c][2.5cm][c]{0.85\linewidth}{\centering\footnotesize\textcolor{midgray}{[Missing image]}}}%
    }%
  }%
}

% --- Section count for navigation (must match actual \section{} count) ---
\setcounter{mlsystotalsections}{8}

\title{Data Engineering}
\author{Vijay Janapa Reddi}
\institute{Harvard University}
\date{}

\begin{document}

% =============================================================================
% TITLE SLIDE
% =============================================================================
\mlsystitle{Data Engineering}{Data Is the Source Code of ML Systems}{cover_data_engineering.png}

% =============================================================================
% LEARNING OBJECTIVES
% =============================================================================
\begin{frame}{Learning Objectives}
\note{[2 min] Walk through objectives. Emphasize that data engineering consumes
60--80\% of ML project effort. Ask: ``Who here has spent more time cleaning
data than training models?''}

\small
\begin{enumerate}
  \item Explain how \textbf{Data Cascades} propagate errors and apply the \textbf{Four Pillars framework}
  \item Evaluate \textbf{data acquisition strategies} using cost--quality trade-offs
  \item Compare \textbf{batch vs.\ streaming} ingestion and \textbf{ETL vs.\ ELT} patterns
  \item Implement \textbf{training-serving consistency} and operationalize \textbf{drift detection}
  \item Build \textbf{data labeling systems} balancing accuracy, throughput, and cost
  \item Evaluate \textbf{storage architectures} and file formats for ML workloads
  \item Identify \textbf{Data Debt} categories and apply systematic debugging
\end{enumerate}

\end{frame}

\begin{frame}{Visual Language}
\note{[1 min] Explain the semantic color system used throughout the course.
These colors are consistent across all diagrams and slides.}

\small
Throughout this course, colors carry meaning:

\vspace{0.3cm}
\begin{columns}[T]
  \begin{column}{0.45\textwidth}
    \begin{mlsyscard}{computestroke}
    \textbf{Blue} --- Compute / Processing\\
    {\footnotesize GPU ops, forward/backward pass, inference}
    \end{mlsyscard}
    \vspace{0.15cm}
    \begin{mlsyscard}{datastroke}
    \textbf{Green} --- Data / Memory\\
    {\footnotesize Data flow, caches, healthy paths}
    \end{mlsyscard}
  \end{column}
  \begin{column}{0.45\textwidth}
    \begin{mlsyscard}{routingstroke}
    \textbf{Orange} --- Routing / Scheduling\\
    {\footnotesize Load balancers, batch windows}
    \end{mlsyscard}
    \vspace{0.15cm}
    \begin{mlsyscard}{errorstroke}
    \textbf{Red} --- Error / Cost / Bottleneck\\
    {\footnotesize Loss, decode phase, waste}
    \end{mlsyscard}
  \end{column}
\end{columns}
\end{frame}


% =============================================================================
\section{Data as Code}
% =============================================================================

\begin{frame}{The Dataset Compilation Metaphor}
\note{[3 min] Core reframe: data engineering is not ``data cleaning'' but
\emph{dataset compilation}. Walk through each analogy pair. Key insight:
datasets must be versioned, tested, debugged --- just like code.
Ask: ``What happens if you delete a training row?''}

% --- Layout: FULL-WIDTH IMAGE ---
\centering
\safeimg[width=0.92\textwidth,height=4.8cm,keepaspectratio]{ch04-dataset-compiler.pdf}

\vspace{0.1cm}
\mlsysconcept{Key insight:}{Deleting a row of training data = deleting a line of code. Retraining = recompiling.}

\end{frame}

\begin{frame}{Data Cascades: Upstream Errors Amplify}
\note{[3 min] Sambasivan et al.\ 2021. Cascade issues take a median of 4 weeks
to discover. By that point, the only fix is often a complete rebuild.
This motivates the Four Pillars framework coming next.
Ask: ``Why are data bugs harder to find than code bugs?''}

\small
\begin{columns}[T]
  \begin{column}{0.55\textwidth}
    \textbf{The failure pattern:}
    \begin{itemize}\setlength\itemsep{0pt}
      \item Bad data at collection $\to$ silent propagation
      \item Errors amplify through every pipeline stage
      \item 4-week median time to discovery
      \item Often requires \textbf{complete system rebuild}
    \end{itemize}

    \vspace{0.15cm}
    \begin{mlsyscard}{errorstroke}
    {\footnotesize \textbf{Example:} A zip code field changed from integer to string. Leading zeros lost. The model rejected all applicants from affected regions --- \emph{silently}.}
    \end{mlsyscard}
  \end{column}
  \begin{column}{0.42\textwidth}
    \vspace{0.3cm}
    \renewcommand{\arraystretch}{1.2}
    {\footnotesize
    \begin{tabular}{@{}ll@{}}
      \toprule
      \textbf{Traditional SW} & \textbf{ML Data} \\
      \midrule
      Compile error    & Silent degradation \\
      Crashes loudly   & Degrades silently \\
      Debug code       & Debug pipeline \\
      \bottomrule
    \end{tabular}
    }
  \end{column}
\end{columns}
\mlsyscite{Sambasivan et al., ``Everyone Wants to Do the Model Work,'' CHI 2021}

\end{frame}

% =============================================================================
\section{Physics of Data}
% =============================================================================

\begin{frame}{The Energy-Movement Invariant}
\note{[3 min] Horowitz 2014. Moving a bit costs 100--1000$\times$ more energy
than computing on it. This is \emph{why} deduplication is the highest-leverage
optimization. Connect to Iron Law: the Data Term dominates energy.
If short on time: focus on the 170$\times$ DRAM vs.\ compute gap.}

% --- Layout: FULL-WIDTH IMAGE ---
\centering
\safeimg[width=0.88\textwidth,height=4.5cm,keepaspectratio]{ch04-energy-movement.pdf}

\vspace{0.15cm}
\mlsysalert{Systems Implication:}{Pruning 50\% of training data eliminates the \emph{most expensive} stages of the pipeline.}

\end{frame}

\begin{frame}{The Feeding Problem: GPU Starvation}
\note{[3 min] An A100 can process 40K images/sec but a standard disk delivers
only 250 MB/s. The GPU sits $>$98\% idle. This is the Iron Law's Data Term
dominating. Adding more GPUs yields zero speedup.
Ask: ``If you have \$10K to spend, should you buy another GPU or faster storage?''}

\small
\begin{columns}[T]
  \begin{column}{0.48\textwidth}
    \safeimg[width=\textwidth,height=5.5cm,keepaspectratio]{ch04-feeding-problem.pdf}
  \end{column}
  \begin{column}{0.48\textwidth}
    \textbf{The mismatch:}

    \vspace{0.15cm}
    \textcolor{computestroke}{\textbf{A100 GPU capacity}}\\
    {\footnotesize 40K images/sec $\to$ needs 23 GB/s}

    \vspace{0.1cm}
    \textcolor{errorstroke}{\textbf{Standard cloud disk}}\\
    {\footnotesize 250 MB/s actual throughput}

    \vspace{0.2cm}
    \begin{mlsyscard}{errorstroke}
    {\footnotesize \textbf{Feeding Tax:} $>$98\% GPU idle time.\\[0.05cm]
    Iron Law: $T_{\text{step}} = \max(T_{\text{compute}}, T_{\text{io}})$\\[0.05cm]
    The \textcolor{datastroke}{Data Term} dominates.}
    \end{mlsyscard}
  \end{column}
\end{columns}

\end{frame}

% --- ACTIVE LEARNING 1: Exercise ---
\begin{frame}{Your Turn: Calculate the Feeding Tax}
\note{[3 min] Give students 90 seconds. Answer: Required BW = 500 MB/s.
Disk provides 100 MB/s. Feeding tax = 80\%. System is I/O bound.
The fix is faster storage, not a faster GPU.
PEER INSTRUCTION PROTOCOL:
1. Present problem (30s). 2. Individual vote (60s).
3. If 30-70\% correct: pair discussion (90s), then re-vote.
4. Reveal and explain.}

\footnotesize
\begin{columns}[T]
  \begin{column}{0.56\textwidth}
    {\small\bfseries Diagnose the I/O Bottleneck}

    \vspace{0.1cm}
    A training pipeline processes:
    \begin{itemize}\setlength\itemsep{0pt}
      \item 2,000 images/sec at 250 KB each
      \item Storage delivers 100 MB/s
    \end{itemize}

    \vspace{0.1cm}
    \textbf{Calculate:} (1) Required bandwidth, (2) Feeding tax, (3) Dominant Iron Law term.

    {\scriptsize\textcolor{midgray}{(90 seconds --- then compare with a neighbor)}}
  \end{column}
  \begin{column}{0.40\textwidth}
    \pause
    \begin{mlsyscard}{datastroke}
    \textbf{Solution:}\\[0.05cm]
    {\footnotesize
    Required: $2000 \times 250\text{KB} = 500$ MB/s\\
    Disk: 100 MB/s\\
    Tax: $1 - 100/500 = $ \textbf{80\%}\\[0.05cm]
    \textcolor{datastroke}{\textbf{Data-bound!}} Fix: faster storage.
    }
    \end{mlsyscard}
  \end{column}
\end{columns}

\end{frame}

% =============================================================================
\section{Four Pillars}
% =============================================================================

% --- ACTIVE LEARNING 2: Predict ---
\begin{frame}{Predict: What Makes a Reliable Data System?}
\note{[2 min] Prediction exercise. Give students 60 seconds to write.
Do NOT reveal the framework yet. This primes the Four Pillars.
Ask 2--3 students to share answers.}

\centering
\vspace{1.2cm}
{\Large\bfseries Think--Write--Share}

\vspace{0.6cm}
{\large A zip code field changes from integer to string.\\
Your ML model silently rejects all applicants from a region.\\[0.3cm]
\alert{What system properties would have prevented this?}}

\vspace{0.6cm}
{\normalsize Write down 3--4 properties. \textcolor{midgray}{(60 seconds)}}

\end{frame}

\begin{frame}{The Four Pillars Framework}
\note{[3 min] Reveal after prediction exercise. Quality, Reliability,
Scalability, Governance. Emphasize interdependence: strengthening one
pillar creates tension with another. This is the organizing framework
for all data engineering decisions.}

% --- Layout: FULL-WIDTH IMAGE ---
\centering
\safeimg[width=0.92\textwidth,height=4.8cm,keepaspectratio]{ch04-four-pillars.pdf}

\vspace{0.1cm}
\mlsysinsight{Framework:}{Every pipeline decision is evaluated against all four pillars simultaneously.}

\end{frame}

\begin{frame}{Four Pillars: Trade-off Tensions}
\note{[2 min] Concrete tension examples. Stricter validation slows throughput
(Quality vs.\ Scalability). More redundancy means more data copies to govern
(Reliability vs.\ Governance). No free lunch in pipeline design.}

\footnotesize
\begin{columns}[T]
  \begin{column}{0.48\textwidth}
    \textbf{Quality $\leftrightarrow$ Scalability}\\[0.1cm]
    {\footnotesize Stricter schema validation catches more errors but reduces pipeline throughput. A 100\% label review catches all errors but costs 10--50$\times$ more.}

    \vspace{0.2cm}
    \textbf{Reliability $\leftrightarrow$ Governance}\\[0.1cm]
    {\footnotesize Redundant data copies improve fault tolerance but multiply the governance surface. Each copy must be versioned, access-controlled, and tracked.}
  \end{column}
  \begin{column}{0.48\textwidth}
    \textbf{Scalability $\leftrightarrow$ Quality}\\[0.1cm]
    {\footnotesize Web scraping scales to billions but requires elaborate cleaning. Curated datasets are high-quality but cap at thousands.}

    \vspace{0.2cm}
    \begin{mlsyscard}{crimson}
    {\footnotesize \textbf{The engineering discipline:} Navigate these tensions systematically, not ad hoc. Every decision shifts the balance.}
    \end{mlsyscard}
  \end{column}
\end{columns}

\end{frame}

% =============================================================================
\section{Data Acquisition}
% =============================================================================

\begin{frame}{Data Acquisition Strategies}
\note{[3 min] Four main strategies, each with different cost--quality trade-offs.
Use the KWS case study: 23 million audio samples, 50 languages.
No single approach works alone.
Ask: ``Which strategy would you use for a medical imaging system?''}

\scriptsize
\renewcommand{\arraystretch}{1.05}
\begin{tabular}{@{}lp{2cm}p{2cm}p{2cm}p{2cm}@{}}
  \toprule
  & \textbf{Curated} & \textbf{Crowdsource} & \textbf{Web Scrape} & \textbf{Synthetic} \\
  \midrule
  \textbf{Scale}   & 1K--10M & 10K--100M & 1M--1B+ & Unlimited \\
  \textbf{Quality} & High & Medium & Low (noisy) & Medium \\
  \textbf{Cost}    & Low (reuse) & Medium & Low (infra) & Medium \\
  \textbf{Risk}    & Overfitting & Annotator bias & Legal, noise & Model collapse \\
  \bottomrule
\end{tabular}

\vspace{0.1cm}
\begin{mlsyscard}{datastroke}
{\scriptsize \textbf{KWS:} Curated Speech Commands (baseline) + crowdsourced MSWC (50 languages) + synthetic TTS (rare accents) + web-scraped noise. No single source suffices.}
\end{mlsyscard}

\end{frame}

\begin{frame}{Labeling: The Serial Bottleneck}
\note{[3 min] Labeling is the most expensive stage. The 1000$\times$ Rule:
labeling costs 1,000--3,000$\times$ more than model training compute.
KWS: labeling 23M samples at 10 sec each = 65,000 hours = 32 person-years.
Ask: ``Why can't we just parallelize labeling infinitely?''}

\small
\begin{columns}[T]
  \begin{column}{0.55\textwidth}
    \textbf{The cost hierarchy:}
    \vspace{0.1cm}
    \renewcommand{\arraystretch}{1.15}
    {\footnotesize
    \begin{tabular}{@{}lr@{}}
      \toprule
      \textbf{Label Type} & \textbf{Cost / Example} \\
      \midrule
      Classification (crowd)     & \$0.01--0.10 \\
      Bounding box               & \$0.10--1.00 \\
      Segmentation mask          & \$1--10 \\
      Expert (medical)           & \$10--100 \\
      \bottomrule
    \end{tabular}
    }

    \vspace{0.15cm}
    \mlsysalert{The 1,000$\times$ Rule:}{Labeling costs 1,000--3,000$\times$ more than model training compute.}
  \end{column}
  \begin{column}{0.42\textwidth}
    \begin{mlsyscard}{routingstroke}
    \textbf{KWS Scale:}\\[0.1cm]
    {\footnotesize
    23M samples $\times$ 10 sec each\\
    = 65,000 hours\\
    = \textbf{32 person-years}\\[0.1cm]
    Budget: \$150K\\
    Timeline: 6 months\\[0.1cm]
    \textbf{Solution:} AI-assisted labeling + consensus for ambiguous cases only.
    }
    \end{mlsyscard}
  \end{column}
\end{columns}

\end{frame}

% =============================================================================
\section{Pipeline Design}
% =============================================================================

\begin{frame}{Pipeline Architecture: ETL vs.\ ELT}
\note{[3 min] Two major patterns. ETL: transform before loading (saves storage,
harder schema changes). ELT: load raw, transform on query (flexible, needs
governance). Most modern ML pipelines use ELT with data lakes.
If short: focus on the trade-off row.}

\footnotesize
\renewcommand{\arraystretch}{1.2}
\begin{tabular}{@{}lll@{}}
  \toprule
  & \textbf{ETL (Extract-Transform-Load)} & \textbf{ELT (Extract-Load-Transform)} \\
  \midrule
  \textbf{Transform}  & Before storage           & After storage (on query) \\
  \textbf{Storage}     & Clean, smaller           & Raw, larger \\
  \textbf{Flexibility} & Schema changes are hard   & Schema-on-read is easy \\
  \textbf{Governance}  & Easier (structured)       & Harder (data swamp risk) \\
  \textbf{ML Use}      & Feature engineering       & Exploratory, training \\
  \textbf{Cost}        & Higher engineering upfront & Higher storage, lower eng. \\
  \bottomrule
\end{tabular}

\vspace{0.15cm}
\begin{mlsyscard}{computestroke}
{\footnotesize \textbf{The trade-off:} ETL reduces storage footprint at the expense of higher engineering overhead during schema changes. ELT preserves raw data for flexibility but risks ``data swamps'' without governance.}
\end{mlsyscard}

\end{frame}

\begin{frame}{Batch vs.\ Streaming Ingestion}
\note{[2 min] Batch: high throughput, periodic updates (training).
Streaming: low latency, continuous updates (serving).
Most ML systems use both: batch for training, streaming for serving features.
Ask: ``When would you choose streaming over batch for training?''}

\small
\begin{columns}[T]
  \begin{column}{0.48\textwidth}
    \begin{mlsyscard}{computestroke}
    \textbf{Batch Ingestion}\\[0.1cm]
    {\footnotesize
    \textbf{When:} Periodic, high-volume\\
    \textbf{Latency:} Minutes to hours\\
    \textbf{Use:} Training data, feature eng.\\
    \textbf{Tools:} Spark, Airflow, dbt\\[0.05cm]
    \textbf{Pro:} Simple, high throughput\\
    \textbf{Con:} Stale data between runs
    }
    \end{mlsyscard}
  \end{column}
  \begin{column}{0.48\textwidth}
    \begin{mlsyscard}{routingstroke}
    \textbf{Stream Ingestion}\\[0.1cm]
    {\footnotesize
    \textbf{When:} Continuous, real-time\\
    \textbf{Latency:} Seconds to minutes\\
    \textbf{Use:} Feature serving, monitoring\\
    \textbf{Tools:} Kafka, Flink, Kinesis\\[0.05cm]
    \textbf{Pro:} Fresh data, low latency\\
    \textbf{Con:} Complex, expensive ops
    }
    \end{mlsyscard}
  \end{column}
\end{columns}

\vspace{0.15cm}
{\small\mlsysconcept{Most ML systems use both:}{Batch for training pipelines, streaming for serving features.}}

\end{frame}

\begin{frame}{Quality Through Validation Gates}
\note{[2 min] Every pipeline stage needs validation. Schema validation catches
structural errors. Statistical monitoring catches distribution drift.
The self-driving car example: 15\% of LiDAR labels were misaligned by 10--20 cm.
Schema checks passed; only statistical monitoring caught it.}

\small
\begin{columns}[T]
  \begin{column}{0.55\textwidth}
    \textbf{Two validation layers:}

    \vspace{0.15cm}
    \textcolor{datastroke}{\textbf{1. Schema Validation}} (\emph{structural})\\
    {\footnotesize Type checks, range constraints, null detection.\\
    Catches: format errors, missing fields.\\
    Misses: semantically wrong but valid data.}

    \vspace{0.15cm}
    \textcolor{computestroke}{\textbf{2. Statistical Monitoring}} (\emph{distributional})\\
    {\footnotesize PSI, KL divergence, feature distributions.\\
    Catches: drift, subtle corruption, bias shifts.\\
    Requires: historical baseline to compare against.}
  \end{column}
  \begin{column}{0.42\textwidth}
    \begin{mlsyscard}{errorstroke}
    \textbf{War Story:}\\[0.1cm]
    {\footnotesize 15\% of LiDAR labels misaligned by 10--20 cm at a self-driving company.\\[0.05cm]
    Every record was \emph{structurally valid}.\\[0.05cm]
    Persisted for 3 months.\\[0.05cm]
    Only statistical monitoring caught it.}
    \end{mlsyscard}
  \end{column}
\end{columns}

\end{frame}

% =============================================================================
\section{Drift \& Skew}
% =============================================================================

\begin{frame}{Three Types of Distribution Drift}
\note{[3 min] Walk through each panel. Covariate: inputs change (detectable).
Label: output frequencies change (detectable). Concept: the mapping changes
(needs ground truth --- hardest to detect). Concept drift is delayed because
labels arrive late. This is the most dangerous silent failure.}

% --- Layout: FULL-WIDTH IMAGE ---
\centering
\safeimg[width=0.92\textwidth,height=4.8cm,keepaspectratio]{ch04-drift-types.pdf}

\vspace{0.1cm}
\mlsysalert{Key:}{Concept drift is the most dangerous --- rules change but detection requires ground truth labels that arrive late.}

\end{frame}

\begin{frame}{Training-Serving Skew}
\note{[3 min] The \#1 cause of ML deployment failures. Different normalization
in training vs.\ serving causes 20--40\% accuracy drop, \emph{silently}.
30--40\% of initial deployments at Uber suffered this.
Solutions: shared code, stored params, feature stores, idempotent transforms.}

% --- Layout: FULL-WIDTH IMAGE ---
\centering
\safeimg[width=0.92\textwidth,height=4.5cm,keepaspectratio]{ch04-training-serving-skew.pdf}

\vspace{0.1cm}
\small
\mlsysconcept{Non-negotiable:}{Any transform applied during training must be applied \emph{identically} during serving.}

\end{frame}

\begin{frame}{Quick Check}
\note{[0.5 min] Micro-retrieval. Give students 30 seconds to recall.
This distributed practice improves retention (Roediger \& Karpicke, 2006).}

\centering
\Large\bfseries Quick check --- no peeking.\\[0.5cm]
\normalsize Name two types of distribution drift and which is hardest to detect.\\[0.3cm]
{\small 30 seconds --- then we continue.}
\end{frame}



% --- ACTIVE LEARNING 3: Discussion ---
\begin{frame}{Discussion: Which Drift Type Is Hardest?}
\note{[3 min] Turn-and-talk. Students discuss in pairs for 90 seconds.
Cold-call 2--3 pairs. Point: concept drift is hardest because you cannot
detect it without ground truth labels, and labels arrive late.
Covariate shift is easiest because you only need to monitor inputs.}

\centering
\vspace{0.8cm}
{\large\bfseries Turn and Talk \textcolor{midgray}{(90 seconds)}}

\vspace{0.6cm}
{\large A fraud detection model worked perfectly at launch.\\
Six months later, loss rates have doubled.\\
No code has changed. No data format errors.\\[0.3cm]
\alert{Which type of drift is this? How would you detect it?}}

\vspace{0.6cm}
\begin{columns}[c]
  \begin{column}{0.30\textwidth}\centering\small Covariate\\ Shift\end{column}
  \begin{column}{0.30\textwidth}\centering\small Label\\ Shift\end{column}
  \begin{column}{0.30\textwidth}\centering\small Concept\\ Drift\end{column}
\end{columns}

\end{frame}

\mlsysfocus{The Degradation Equation}{%
$A(t) = A_0 - \alpha \cdot \Delta(P_{\text{train}},\; P_{\text{live}}(t))$\\[0.3cm]
\normalsize Drift detection (PSI, KL divergence) operationalizes this equation\\
into monitoring infrastructure%
}

% =============================================================================
\section{Storage \& Debt}
% =============================================================================

\begin{frame}{ML Storage Hierarchy}
\note{[3 min] Each tier optimizes for a different access pattern. The 50$\times$
speed gap between NVMe and S3 determines daily vs.\ weekly iteration cycles.
KWS: 736 GB on S3 costs \$17/mo, on NVMe costs \$74--221/mo but loads in
2.5 min vs.\ 2 hours. The throughput gap is the binding constraint on
iteration velocity.}

% --- Layout: FULL-WIDTH IMAGE ---
\centering
\safeimg[width=0.88\textwidth,height=4.5cm,keepaspectratio]{ch04-storage-hierarchy.pdf}

\vspace{0.1cm}
\mlsysconcept{Iron Law connection:}{Storage hierarchy directly sets $D_{\text{vol}}/BW$ --- the Data Term.}

\end{frame}

\begin{frame}{Storage Systems: Database vs.\ Warehouse vs.\ Lake}
\note{[2 min] Three architectures for different access patterns. Databases:
high IOPS for serving. Warehouses: high throughput for analytics. Lakes:
flexible for raw data. Most ML orgs use all three, orchestrated through
unified catalogs. If short: just show the table.}

\scriptsize
\renewcommand{\arraystretch}{1.1}
\begin{tabular}{@{}lp{2.5cm}p{2.5cm}p{2.5cm}@{}}
  \toprule
  & \textbf{Database (OLTP)} & \textbf{Warehouse (OLAP)} & \textbf{Data Lake} \\
  \midrule
  \textbf{Optimized for} & High IOPS, low latency & High throughput scans & Capacity, flexibility \\
  \textbf{Data type}     & Structured              & Structured            & Any (schema-on-read) \\
  \textbf{ML use}        & Feature serving         & Feature engineering   & Training, raw storage \\
  \textbf{Access}        & Point lookups (ms)      & Columnar scans (s)    & Sequential reads \\
  \textbf{Risk}          & Scaling limits           & Schema rigidity       & ``Data swamp'' \\
  \bottomrule
\end{tabular}

\vspace{0.15cm}
\begin{mlsyscard}{routingstroke}
{\scriptsize \textbf{Mature ML organizations use all three:} Databases for real-time serving, warehouses for curated analytics, and data lakes for petabyte-scale raw data. The wrong system for a workload creates order-of-magnitude performance penalties.}
\end{mlsyscard}

\end{frame}

\begin{frame}{Data Debt: The Hidden Iceberg}
\note{[3 min] 65--95\% of ML failures trace to data, not models. Four categories
of hidden debt: Freshness, Quality, Schema, Documentation. Unlike code debt
(slower development), data debt causes lower accuracy even when code is perfect.
Data debt compounds silently.}

% --- Layout: FULL-WIDTH IMAGE ---
\centering
\safeimg[width=0.88\textwidth,height=4.5cm,keepaspectratio]{ch04-data-debt-categories.pdf}

\vspace{0.1cm}
\mlsysalert{Key:}{Unlike code debt (slower development), data debt causes \emph{lower accuracy} even when code is perfect.}

\end{frame}

% =============================================================================
\section{Wrap-Up}
% =============================================================================

\begin{frame}{Fallacies}
\note{[2 min] Three common misconceptions with quantitative evidence.}

\small
\textbf{Fallacy:} \textit{More data always improves model performance.}\\
{\footnotesize Beyond a threshold, test loss follows a power law: 10$\times$ more data often reduces error by $<$1 percentage point, while costs scale linearly. Smart selection outperforms naive accumulation.}

\vspace{0.15cm}
\textbf{Fallacy:} \textit{High training accuracy indicates production readiness.}\\
{\footnotesize Training accuracy measures fit to historical data. Training-serving skew, distribution drift, and coverage gaps cause 99\% validation accuracy to fail in deployment.}

\vspace{0.15cm}
\textbf{Fallacy:} \textit{Synthetic data can fully replace real-world collection.}\\
{\footnotesize Synthetic data inherits generator biases. A KWS system trained on synthesized speech fails on real accents and background noise. Combine real data for coverage with synthetic for scale.}

\end{frame}

\begin{frame}{Pitfalls}
\note{[2 min] Three operational pitfalls.}

\footnotesize
\textbf{Pitfall:} \textit{Treating data preprocessing as a one-time task.}\\
Data distributions drift continuously. A pipeline validated at launch degrades silently. Production systems require continuous monitoring (PSI, KL divergence) and automated retraining triggers.

\vspace{0.15cm}
\textbf{Pitfall:} \textit{Ignoring training-serving skew until deployment.}\\
Feature computation differences are the leading cause of ML deployment failures. Feature stores and consistency contracts should be architectural requirements from day one.

\vspace{0.15cm}
\textbf{Pitfall:} \textit{Neglecting data versioning until model debugging requires it.}\\
Without versioned snapshots, teams cannot determine whether regression stems from labeling changes, schema errors, or genuine drift. Organizations that defer versioning report 2--4$\times$ longer debugging cycles.

\end{frame}


\begin{frame}{Muddiest Point}
\note{[1 min] One-minute paper. Collect responses (physical or digital).
Use responses to open the NEXT lecture with targeted clarification.
Angelo \& Cross (1993) --- powerful diagnostic tool.}

\centering
\Large\bfseries One-minute paper\\[0.5cm]
\normalsize Write down the \textbf{muddiest point} from today's lecture.\\[0.2cm]
{\small What concept was most confusing or unclear?}\\[0.5cm]
{\footnotesize\textcolor{midgray}{(Hand in on your way out --- or submit digitally)}}
\end{frame}

% --- RETRIEVAL PRACTICE ---
\begin{frame}{What Were the Key Ideas?}
\note{[2 min] Retrieval practice. Students write 90 seconds, no notes.
Do NOT show next slide yet. Walk around the room.}

\centering
\vspace{1.5cm}
{\Large\bfseries Close your notes.}

\vspace{0.8cm}
{\large Write down the \textbf{4 most important concepts} from today.}

\vspace{0.8cm}
{\normalsize\textcolor{midgray}{90 seconds --- no peeking.}}

\end{frame}

\begin{frame}{Key Takeaways}
\note{[2 min] Reveal. Walk through each bullet. Emphasize quantitative
anchors: 60--80\% effort, $>$98\% feeding tax, 1,000$\times$ labeling cost,
65--95\% data failures, 50$\times$ storage speed gap.}

\scriptsize
\begin{itemize}\setlength\itemsep{0pt}
  \item \textbf{Data is source code}: Datasets must be versioned, tested, debugged. Deleting a row = deleting a line of code.
  \item \textbf{Data cascades}: Upstream errors amplify through every stage. Four Pillars (Quality, Reliability, Scalability, Governance) organize prevention.
  \item \textbf{Energy-movement invariant}: Moving a bit costs 100--1,000$\times$ more than computing. Deduplication is the highest-leverage optimization.
  \item \textbf{Training-serving consistency}: Non-negotiable. Feature stores and shared code prevent the \#1 deployment failure.
  \item \textbf{Drift detection}: Degradation Equation ($A(t) = A_0 - \alpha \cdot \Delta$) operationalized through PSI and KL monitoring.
  \item \textbf{Storage hierarchy}: 50$\times$ NVMe vs.\ S3 speed gap determines daily vs.\ weekly iteration. Choose by $D_{\text{vol}}/BW$.
  \item \textbf{Data debt compounds}: 65--95\% of ML failures trace to data. Allocate sustained engineering capacity.
\end{itemize}

\end{frame}

\begin{frame}{References}
\note{[1 min] Point students to canonical papers.}

\small
\mlsysref{Sambasivan+21}{Sambasivan et al. ``Everyone Wants to Do the Model Work, Not the Data Work.'' CHI 2021.}
\mlsysref{Sculley+15}{Sculley et al. ``Hidden Technical Debt in ML Systems.'' NeurIPS 2015.}
\mlsysref{Horowitz14}{M. Horowitz. ``Computing's Energy Problem.'' ISSCC 2014.}
\mlsysref{Hestness+17}{Hestness et al. ``Deep Learning Scaling is Predictable, Empirically.'' 2017.}
\mlsysref{Polyzotis+18}{Polyzotis et al. ``Data Lifecycle Challenges in Production ML.'' SIGMOD 2018.}
\mlsysref{Northcutt+21}{Northcutt et al. ``Pervasive Label Errors in Test Sets.'' NeurIPS 2021.}

\end{frame}

\begin{frame}{Next Lecture: Neural Network Computation}
\note{[1 min] Forward hook. The dataset compiler has produced its output.
A compiled binary does nothing until it runs on hardware. We turn next to
the mathematical foundations of learning --- neurons, weights, gradients.}

\centering
\vspace{0.5cm}
{\large The dataset compiler has produced its output:}\\[0.1cm]
{\large a clean, versioned, optimized training set.}

\vspace{0.5cm}
{\large A compiled binary does nothing until it runs on hardware.}

\vspace{0.5cm}

\begin{mlsyscard}{crimson}
\centering
{\large\bfseries What happens when data meets computation?}\\[0.2cm]
{\normalsize Neurons, weights, gradients, and the math of learning.}
\end{mlsyscard}

\vspace{0.3cm}
{\small\textbf{Next:} Ch5 --- Neural Network Computation}

\end{frame}


% =============================================================================
% BACKUP SLIDES
% =============================================================================
\appendix

\begin{frame}{Backup: PSI Calculation Walkthrough}
\note{Use for additional depth or if students need alternative explanation.}
\small
Use if students want a worked PSI example.\\small PSI $= \sum (p_i - q_i) \times \ln(p_i / q_i)$ where $p_i$ = actual proportion, $q_i$ = expected proportion in bin $i$.\\ PSI $< 0.1$: no significant shift.\\ PSI $0.1$--$0.2$: moderate shift, investigate.\\ PSI $> 0.2$: significant shift, retrain.
\end{frame}

\begin{frame}{Backup: File Format Comparison for ML}
\note{Use for additional depth or if students need alternative explanation.}
\small
Use if students ask about Parquet vs.\\ CSV vs.\\ TFRecord.\\small Parquet: columnar, compressed, 10$\times$ smaller than CSV, fast column scans.\\ TFRecord: sequential, optimized for TensorFlow pipelines.\\ CSV: human-readable but 10$\times$ larger and 5$\times$ slower to load.
\end{frame}

\end{document}
